\documentclass[../../main.tex]{subfiles}% 注意这里的文件路径不能用 ./main.tex ,否则用latexmk编译子文件会报错
\graphicspath{{\subfix{./image/}}} % 指定图片目录,后续可以直接使用图片文件名
% 注意这里的文件路径不能用 ../../image/ ,否则用latexmk编译子文件会报错

% 例如:
% \begin{figure}[H]
% \centering
% \includegraphics[scale=0.3]{图.png}
% \caption{}
% \label{figure:图}
% \end{figure}
% 注意:上述\label{}一定要放在\caption{}之后,否则引用图片序号会只会显示??.

\begin{document}

\section{重积分换元法}

\begin{theorem}[重积分换元法]\label{theorem:重积分换元法}
\begin{enumerate}
\item 考虑换元 $w:\begin{cases}x = x(u,v)\\y = y(u,v)\end{cases}$ ,则有
\begin{align*}
\iint_D f(x,y) \mathrm{d}x\mathrm{d}y= \iint_{w^{-1}(D)} f(x(u,v),y(u,v)) \cdot |\det J| dudv,
\end{align*}
这里
\begin{align*}
J \triangleq \frac{\partial(x,y)}{\partial(u,v)} = {\textstyle \begin{pmatrix}\frac{\partial x}{\partial u}&\frac{\partial x}{\partial v}\\\frac{\partial y}{\partial u}&\frac{\partial y}{\partial v}\end{pmatrix}}.
\end{align*}
\item 考虑换元 $w:\begin{cases}x = x(u,v,w)\\y = y(u,v,w)\\z = z(u,v,w)\end{cases}$ ,则有
\begin{align*}
\iiint_D f(x,y,z) \mathrm{d}x\mathrm{d}y\mathrm{d}z = \iiint_{w^{-1}(D)} f(x(u,v,w),y(u,v,w),z(u,v,w)) \cdot |\det J| \mathrm{d}u\mathrm{d}v\mathrm{d}w,
\end{align*}
这里
\begin{align*}
J \triangleq \frac{\partial(x,y,z)}{\partial(u,v,w)} = \begin{pmatrix}\frac{\partial x}{\partial u}&\frac{\partial x}{\partial v}&\frac{\partial x}{\partial w}\\\frac{\partial y}{\partial u}&\frac{\partial y}{\partial v}&\frac{\partial y}{\partial w}\\\frac{\partial z}{\partial u}&\frac{\partial z}{\partial v}&\frac{\partial z}{\partial w}\end{pmatrix}.
\end{align*}
\end{enumerate}
\end{theorem}
\begin{remark}
\begin{enumerate}[(1)]
\item 记忆结论
\begin{align*}
\left[\frac{\partial(u,v)}{\partial(x,y)}\right]^{-1} = \frac{\partial(x,y)}{\partial(u,v)}, \quad \left[\frac{\partial(u,v,w)}{\partial(x,y,z)}\right]^{-1} = \frac{\partial(x,y,z)}{\partial(u,v,w)},
\end{align*}
于是有
\begin{align*}
\left|\det \frac{\partial(u,v)}{\partial(x,y)}\right| \cdot \left|\det \frac{\partial(x,y)}{\partial(u,v)}\right| = 1,\quad  \left|\det \frac{\partial(u,v,w)}{\partial(x,y,z)}\right| \cdot \left|\det \frac{\partial(x,y,z)}{\partial(u,v,w)}\right| = 1.
\end{align*}

\item 换元法的区域确定重点是:\textbf{边界对应边界}.

\item 极坐标,柱坐标,球坐标等方法是换元法的特例,我们课本上记忆的经典确定上下限的几何直观是他们独有的. 对于一般没有,或者很难有几何意义的换元法,只能把换元后的新变量用直角坐标的方法确定上下限. 特别的,极坐标,柱坐标,球坐标等方法也可以视换元后的新变量用直角坐标的方法确定上下限.

\item 计算积分时常用的轮换变换、反射(对称)变换、平移变换都可视作重积分换元,只不过它们的雅克比行列式的绝对值都是1.
\end{enumerate}
\end{remark}

\begin{example}
设$D=\left\{(x,y)\in\mathbb{R}^2:2\leqslant \frac{x}{x^2+y^2},\frac{y}{x^2+y^2}\leqslant4\right\}$，计算$\iint_{D}\frac{1}{xy}\mathrm{d}x\mathrm{d}y$.
\end{example}

\begin{proof}

做换元
\[
\begin{cases}
u=\displaystyle\frac{x}{x^2+y^2}\\[4pt]
v=\displaystyle\frac{y}{x^2+y^2}
\end{cases},
\]
则
\begin{align*}
\left|\det \frac{\partial(u,v)}{\partial(x,y)}\right|
=\left|\det \begin{pmatrix}
\displaystyle\frac{y^2-x^2}{(x^2+y^2)^2} & \displaystyle-\frac{2xy}{(x^2+y^2)^2}\\[6pt]
\displaystyle-\frac{2xy}{(x^2+y^2)^2} & \displaystyle\frac{x^2-y^2}{(x^2+y^2)^2}
\end{pmatrix}\right|
=\frac{1}{(x^2+y^2)^2}.
\end{align*}
于是我们有
\begin{align*}
\iint_{D}\frac{1}{xy}\mathrm{d}x\mathrm{d}y
=\int_{2}^{4}\int_{2}^{4}\frac{(x^2+y^2)^2}{xy}\mathrm{d}u\mathrm{d}v
=\int_{2}^{4}\int_{2}^{4}\frac{1}{uv}\mathrm{d}u\mathrm{d}v=\ln^2 2.
\end{align*}
\end{proof}

\begin{example}
设$\Omega$由$z=y^2$，$z=4y^2$，$z=x$，$z=2x$，$z=2$围成的$y>0$的区域，计算
\begin{align*}
\iiint_{\Omega}\frac{z\sqrt{z}}{y^3}\cos\frac{z}{y^2}\mathrm{d}x\mathrm{d}y\mathrm{d}z.
\end{align*}
\end{example}

\begin{proof}

做换元
\[
\begin{cases}
u=\displaystyle\frac{z}{y^2}\\[4pt]
v=\displaystyle\frac{z}{x}\\[4pt]
w=z
\end{cases}
\iff
\begin{cases}
x=\displaystyle\frac{z}{v}\\[4pt]
y=\displaystyle\sqrt{\frac{z}{u}}\\[4pt]
z=w
\end{cases},
\]
则计算得
\begin{align*}
\left|\det \frac{\partial(x,y,z)}{\partial(u,v,w)}\right|
=\left|\det \begin{pmatrix}
0 & -\displaystyle\frac{w}{v^2} & \displaystyle\frac{1}{v}\\[6pt]
-\displaystyle\frac12\sqrt{\frac{w}{u^3}} & 0 & \displaystyle\frac{1}{2\sqrt{uw}}\\[6pt]
0 & 0 & 1
\end{pmatrix}\right|
=\frac{1}{2v^2}\left(\frac{w}{u}\right)^{\frac32}.
\end{align*}
于是我们有
\begin{align*}
\iiint_{\Omega}\frac{z\sqrt{z}}{y^3}\cos\frac{z}{y^2}\mathrm{d}x\mathrm{d}y\mathrm{d}z
=\frac12\int_{1}^{4}\cos u\mathrm{d}u\int_{1}^{2}\frac{1}{v^2}\mathrm{d}v\int_{0}^{2}w^{\frac32}\mathrm{d}w
=\frac{2\sqrt{2}}{5}(\sin4-\sin1).
\end{align*}
\end{proof}

\begin{theorem}[第一类曲线/曲面积分换元法]\label{theorem:第一类曲线曲面积分换元法}
\begin{enumerate}[(1)]
\item 我们约定
\begin{align*}
D_F \triangleq \begin{pmatrix} \dfrac{\partial F}{\partial x}\\[4pt] \dfrac{\partial F}{\partial y} \end{pmatrix},\quad \|D_F\| \triangleq \sqrt{\left|\frac{\partial F}{\partial x}\right|^2+\left|\frac{\partial F}{\partial y}\right|^2}.
\end{align*}
于是做换元
\[
\begin{cases}
x=x(u,v)\\
y=y(u,v)
\end{cases},
\]
我们有
\begin{align*}
\int_{F(x,y)=0} f(x,y)\mathrm{d}s
=\int_{F(x(u,v),y(u,v))=0} f(x(u,v),y(u,v))\frac{|\det J|\cdot\|D_F\|}{\|J^\top\cdot D_F\|}\mathrm{d}s,
\end{align*}
这里
\begin{align*}
J \triangleq \frac{\partial(x,y,z)}{\partial(u,v,w)} = \begin{pmatrix}\frac{\partial x}{\partial u}&\frac{\partial x}{\partial v}&\frac{\partial x}{\partial w}\\\frac{\partial y}{\partial u}&\frac{\partial y}{\partial v}&\frac{\partial y}{\partial w}\\\frac{\partial z}{\partial u}&\frac{\partial z}{\partial v}&\frac{\partial z}{\partial w}\end{pmatrix}.
\end{align*}

\item 我们约定
\begin{align*}
D_F\triangleq \begin{pmatrix} \dfrac{\partial F}{\partial x}\\[4pt] \dfrac{\partial F}{\partial y}\\[4pt] \dfrac{\partial F}{\partial z} \end{pmatrix},\quad \|D_F\| \triangleq \sqrt{\left|\frac{\partial F}{\partial x}\right|^2+\left|\frac{\partial F}{\partial y}\right|^2+\left|\frac{\partial F}{\partial z}\right|^2}.
\end{align*}
于是做换元
\[
\begin{cases}
x=x(u,v,w)\\
y=y(u,v,w)\\
z=z(u,v,w)
\end{cases},
\]
我们有
\begin{align*}
\iint_{F(x,y,z)=0} f(x,y,z)\mathrm{d}S
=\iint_{F(x(u,v,w),y(u,v,w),z(u,v,w))=0} f(x(u,v,w),y(u,v,w),z(u,v,w))\frac{|\det J|\cdot\|D_F\|}{\|J^\top\cdot D_F\|}\mathrm{d}S,
\end{align*}
这里
\begin{align*}
J \triangleq \frac{\partial(x,y,z)}{\partial(u,v,w)} = \begin{pmatrix}\frac{\partial x}{\partial u}&\frac{\partial x}{\partial v}&\frac{\partial x}{\partial w}\\\frac{\partial y}{\partial u}&\frac{\partial y}{\partial v}&\frac{\partial y}{\partial w}\\\frac{\partial z}{\partial u}&\frac{\partial z}{\partial v}&\frac{\partial z}{\partial w}\end{pmatrix}.
\end{align*}
\end{enumerate}
\end{theorem}

\begin{remark}
\begin{enumerate}[(1)]
\item 设
\[
\begin{cases}
x=a+r_1 u\\
y=b+r_2 v
\end{cases},
\]
可以看到
\begin{align*}
\frac{|\det J|\cdot\|D_F\|}{\|J^\top\cdot D_F\|}=\frac{r_1 r_2\sqrt{\left|\frac{\partial F}{\partial x}\right|^2+\left|\frac{\partial F}{\partial y}\right|^2}}{\sqrt{r_1^2\left|\frac{\partial F}{\partial x}\right|^2+r_2^2\left|\frac{\partial F}{\partial y}\right|^2}},
\end{align*}
因此这种换元是十分复杂的除非$r_1=r_2=r$。当$r_1=r_2=r$，我们有
\begin{align*}
\frac{|\det J|\cdot\|D_F\|}{\|J^\top\cdot D_F\|}=r.
\end{align*}
此外注意到换元$\begin{cases}x=a+ru\\y=b+rv\end{cases}$对二重积分而言是$r^2$，但对曲线积分却是$r$。类似的对三重积分是$r^3$，但对曲面积分却是$r^2$。一般的，对于$\mathbb{R}^n$中的曲面积分，其$r$的次方要比重积分少$1$，这也反映了曲面积分是低一维的积分。

\item 重积分正交变换下显然不变，因为换元得到的雅克比矩阵是正交矩阵，而正交矩阵行列式为$\pm1$。

\item 曲线曲面积分正交变换下也不变，因为设$\begin{pmatrix}x\\y\end{pmatrix}=J^\top\begin{pmatrix}u\\v\end{pmatrix}$，$J$是正交矩阵，则
\begin{align*}
\|J^\top\cdot D_F\|=\sqrt{D_F^\top JJ^\top D_F}=\|D_F\| \implies \frac{|\det J|\cdot\|D_F\|}{\|J^\top\cdot D_F\|}=1.
\end{align*}
\end{enumerate}
\end{remark}

\begin{example}
计算$\iint_{\Omega}(x+y+z)\mathrm{d}S$，这里$\Omega=\{(x,y,z)\in\mathbb{R}^3:x^2+y^2+z^2=1,x+y+z\geqslant0\}$.
\end{example}

\begin{proof}

待定换元$\begin{pmatrix}x\\y\\z\end{pmatrix}=T\begin{pmatrix}u\\v\\w\end{pmatrix}$，$T$是三阶正交矩阵。可让
\[
\begin{pmatrix}u\\v\\w\end{pmatrix}=T^\top\begin{pmatrix}x\\y\\z\end{pmatrix}=
\begin{pmatrix}
* & * & *\\
* & * & *\\
\frac1{\sqrt3} & \frac1{\sqrt3} & \frac1{\sqrt3}
\end{pmatrix}\begin{pmatrix}x\\y\\z\end{pmatrix}.
\]
于是
\begin{align*}
\iint_{\Omega}(x+y+z)\mathrm{d}S=\sqrt3\iint_{\Omega'} w\mathrm{d}S,\quad \Omega'=\{(u,v,w)\in\mathbb{R}^3:u^2+v^2+w^2=1,w\geqslant0\},
\end{align*}
这里用到了
\[
u^2+v^2+w^2=\begin{pmatrix}u & v & w\end{pmatrix}I_3\begin{pmatrix}u\\v\\w\end{pmatrix}=\begin{pmatrix}x & y & z\end{pmatrix}TI_3T^\top\begin{pmatrix}x\\y\\z\end{pmatrix}=x^2+y^2+z^2=1.
\]
现在我们就知道
\begin{align*}
\iint_{\Omega}(x+y+z)\mathrm{d}S=\sqrt3\iint_{\Omega'} w\mathrm{d}S\xlongequal{\text{投影法}}\sqrt3\iint_{x^2+y^2\leqslant1}1\mathrm{d}x\mathrm{d}y=\sqrt3\pi.
\end{align*}
\end{proof}

\begin{theorem}[第二类曲线积分换元法]\label{theorem:第二类曲线积分换元法}
考虑换元$\begin{cases}x=x(u,v)\\y=y(u,v)\end{cases}$，则
\begin{align*}
\int_L f(x,y)\mathrm{d}x+g(x,y)\mathrm{d}y
&=\int_{L'} f(x(u,v),y(u,v))\mathrm{d}x(u,v)+g(x(u,v),y(u,v))\mathrm{d}y(u,v)\\
&=\int_{L'} f(x(u,v),y(u,v))\left[\frac{\partial x}{\partial u}\mathrm{d}u+\frac{\partial x}{\partial v}\mathrm{d}v\right]+g(x(u,v),y(u,v))\left[\frac{\partial y}{\partial u}\mathrm{d}u+\frac{\partial y}{\partial v}\mathrm{d}v\right],
\end{align*}
这里$L'$是$L$在换元下对应的像，方向如下确定：
\[
\begin{cases}
\dfrac{\partial(x,y)}{\partial(u,v)}>0,\quad \text{与}L\text{同向}\\[4pt]
\dfrac{\partial(x,y)}{\partial(u,v)}<0,\quad \text{与}L\text{反向}.
\end{cases}
\]
\end{theorem}
\begin{remark}
\begin{enumerate}[(1)]
\item 不给出第二类曲面积分换元法，因为从未用上。
\item 第二类曲线积分换元法只需记忆一个方向即可，换元只需要自然的计算。
\end{enumerate}
\end{remark}

\begin{example}
计算
\[
\lim_{r\to+\infty}\oint_{x^2+xy+y^2=r^2}\frac{y\mathrm{d}x-x\mathrm{d}y}{(x^2+y^2)^\alpha},
\]
这里积分方向是逆时针。
\end{example}
\begin{proof}

注意到
\[
x^2+xy+y^2=\begin{pmatrix}x & y\end{pmatrix}\begin{pmatrix}1 & \frac12\\[4pt]\frac12 & 1\end{pmatrix}\begin{pmatrix}x\\y\end{pmatrix},
\]
于是可以求一个正交矩阵$T\in\mathbb{R}^{2\times2}$使得$\begin{pmatrix}1 & \frac12\\\frac12 & 1\end{pmatrix}$对角化。线性代数基本知识告诉我们可取$T=\begin{pmatrix}\frac1{\sqrt2} & -\frac1{\sqrt2}\\[4pt]\frac1{\sqrt2} & \frac1{\sqrt2}\end{pmatrix}$，此时
\[
T^{-1}\begin{pmatrix}1 & \frac12\\\frac12 & 1\end{pmatrix}T=\begin{pmatrix}\frac32 & 0\\0 & \frac12\end{pmatrix}.
\]
做换元$\begin{pmatrix}x\\y\end{pmatrix}=T\begin{pmatrix}u\\v\end{pmatrix}$，即
\[
\begin{cases}
x=\dfrac{u}{\sqrt2}-\dfrac{v}{\sqrt2}\\[4pt]
y=\dfrac{u}{\sqrt2}+\dfrac{v}{\sqrt2}.
\end{cases}
\]
我们可以算得$\det \dfrac{\partial(x,y)}{\partial(u,v)}=1>0$。现在积分方向仍然是逆时针，于是就有
\begin{align*}
\lim_{r\to+\infty}\oint_{x^2+xy+y^2=r^2}\frac{y\mathrm{d}x-x\mathrm{d}y}{(x^2+y^2)^\alpha}
=\lim_{r\to+\infty}\oint_{\frac32u^2+\frac12v^2=r^2}\frac{\frac{u+v}{\sqrt2}\mathrm{d}\frac{u-v}{\sqrt2}-\frac{u-v}{\sqrt2}\mathrm{d}\frac{u+v}{\sqrt2}}{\left(\left(\frac{u-v}{\sqrt2}\right)^2+\left(\frac{u+v}{\sqrt2}\right)^2\right)^\alpha}.
\end{align*}
再做换元$\begin{cases}u=\sqrt3 r\cos\theta\\v=\sqrt2 r\sin\theta\end{cases}$，
\[
=\lim_{r\to+\infty}\oint_{\frac32u^2+\frac12v^2=r^2}\frac{v\mathrm{d}u-u\mathrm{d}v}{(u^2+v^2)^\alpha}
=\lim_{r\to+\infty}\frac{-2r^{2(1-\alpha)}}{\sqrt3}\int_{0}^{2\pi}\frac{\mathrm{d}\theta}{\left(\frac23\cos^2\theta+2\sin^2\theta\right)^\alpha}.
\]
当$\alpha<1$极限不存在，当$\alpha>1$，我们就有
\[
\lim_{r\to+\infty}\oint_{x^2+xy+y^2=r^2}\frac{y\mathrm{d}x-x\mathrm{d}y}{(x^2+y^2)^\alpha}=0.
\]
当$\alpha=1$，我们就有
\begin{align*}
\lim_{r\to+\infty}\oint_{x^2+xy+y^2=r^2}\frac{y\mathrm{d}x-x\mathrm{d}y}{x^2+y^2}
&=\frac{-2}{\sqrt3}\int_{0}^{2\pi}\frac{\mathrm{d}\theta}{\frac23\cos^2\theta+2\sin^2\theta}
=\frac{-2}{\sqrt3}\int_{-\pi}^{\pi}\frac{\mathrm{d}\theta}{\frac23\cos^2\theta+2\sin^2\theta}
=\frac{-8}{\sqrt3}\int_{0}^{\frac{\pi}{2}}\frac{\mathrm{d}\theta}{\frac23\cos^2\theta+2\sin^2\theta}\\
&=\frac{-8}{\sqrt3}\int_{0}^{\frac{\pi}{2}}\frac{\frac1{\cos^2\theta}\mathrm{d}\theta}{\frac23+2\tan^2\theta}
=\frac{-8}{\sqrt3}\int_{0}^{+\infty}\frac{\mathrm{d}\tan\theta}{\frac23+2t^2}
=-2\pi.
\end{align*}
\end{proof}

\begin{theorem}[第一类曲面积分的参数方程法]\label{theorem:第一类曲面积分的参数方程法}
若曲面$F(x,y,z)=0$有参数方程表示$(x(u,v),y(u,v),z(u,v))$。我们记
\[
\Delta(u,v)=\sqrt{\left|\det \frac{\partial(x,y)}{\partial(u,v)}\right|^2+\left|\det \frac{\partial(x,z)}{\partial(u,v)}\right|^2+\left|\det \frac{\partial(y,z)}{\partial(u,v)}\right|^2},
\]
则
\begin{align*}
\iint_{F(x,y,z)=0} f(x,y,z)\mathrm{d}S=\iint f(x(u,v),y(u,v),z(u,v))\cdot \Delta(u,v)\,\mathrm{d}u\mathrm{d}v,
\end{align*}
这里第二个积分的积分区域是参数$(u,v)$的取值区域。
\end{theorem}
\begin{remark}
特别的，对于球坐标换元
\[
\begin{cases}
x=x_0+r\sin\varphi\cos\theta\\
y=y_0+r\sin\varphi\sin\theta\\
z=z_0+r\cos\varphi
\end{cases},\quad \varphi\in[0,\pi],\theta\in[0,2\pi],
\]
我们有$\Delta(\varphi,\theta)=r^2\sin\varphi$.
\end{remark}

\begin{theorem}[余面积公式]\label{theorem:余面积公式}
\begin{enumerate}[(1)]
\item
\begin{align*}
\iint_{r_1\leqslant F(x,y)\leqslant r_2} f(x,y)\mathrm{d}x\mathrm{d}y = \int_{r_1}^{r_2}\mathrm{d}r \int_{F(x,y)=r} \frac{f(x,y)}{\sqrt{\left|\frac{\partial F}{\partial x}\right|^2+\left|\frac{\partial F}{\partial y}\right|^2}}\mathrm{d}S.
\end{align*}
\item
\begin{align*}
\iiint_{r_1\leqslant F(x,y,z)\leqslant r_2} f(x,y,z)\mathrm{d}x\mathrm{d}y\mathrm{d}z = \int_{r_1}^{r_2}\mathrm{d}r \iint_{F(x,y,z)=r} \frac{f(x,y,z)}{\sqrt{\left|\frac{\partial F}{\partial x}\right|^2+\left|\frac{\partial F}{\partial y}\right|^2+\left|\frac{\partial F}{\partial z}\right|^2}}\mathrm{d}S.
\end{align*}
\end{enumerate}
\end{theorem}
\begin{remark}
特别的取$F(x,y)=\sqrt{x^2+y^2}$，$F(x,y,z)=\sqrt{x^2+y^2+z^2}$，我们有
\begin{enumerate}[(1)]
\item
\begin{align*}
\iint_{r_1\leqslant \sqrt{x^2+y^2}\leqslant r_2} f(x,y)\mathrm{d}x\mathrm{d}y = \int_{r_1}^{r_2}\mathrm{d}r \int_{\sqrt{x^2+y^2}=r} f(x,y)\mathrm{d}S.
\end{align*}
\item
\begin{align*}
\iiint_{r_1\leqslant \sqrt{x^2+y^2+z^2}\leqslant r_2} f(x,y,z)\mathrm{d}x\mathrm{d}y\mathrm{d}z = \int_{r_1}^{r_2}\mathrm{d}r \iint_{\sqrt{x^2+y^2+z^2}=r} f(x,y,z)\mathrm{d}S.
\end{align*}
\end{enumerate}
\end{remark}

\begin{example}
对$a,b,c>0$，计算
\begin{align*}
\iint_{\frac{x^2}{a^2}+\frac{y^2}{b^2}+\frac{z^2}{c^2}=1} \frac{1}{\sqrt{\frac{x^2}{a^4}+\frac{y^2}{b^4}+\frac{z^2}{c^4}}}\mathrm{d}S.
\end{align*}
\end{example}
\begin{proof}

利用\hyperref[theorem:余面积公式]{余面积公式}，我们有
\begin{align*}
\int_{0}^{r}\iint_{\frac{x^2}{a^2}+\frac{y^2}{b^2}+\frac{z^2}{c^2}=t}\frac{1}{\sqrt{\frac{x^2}{a^4}+\frac{y^2}{b^4}+\frac{z^2}{c^4}}}\mathrm{d}S\mathrm{d}t
=2\iiint_{\frac{x^2}{a^2}+\frac{y^2}{b^2}+\frac{z^2}{c^2}\leqslant r}1\mathrm{d}x\mathrm{d}y\mathrm{d}z
=2abc\iiint_{x^2+y^2+z^2\leqslant r}1\mathrm{d}x\mathrm{d}y\mathrm{d}z
=\frac{8\pi abc}{3}r^\frac{3}{2}.
\end{align*}
于是
\begin{align*}
\iint_{\frac{x^2}{a^2}+\frac{y^2}{b^2}+\frac{z^2}{c^2}=r}\frac{1}{\sqrt{\frac{x^2}{a^4}+\frac{y^2}{b^4}+\frac{z^2}{c^4}}}\mathrm{d}S
=\frac{8\pi abc}{3}\frac{\mathrm{d}}{\mathrm{d}r}r^\frac{3}{2}=4\pi abc\sqrt{r},\ r>0.
\end{align*}
让$r=1$，即可得
\begin{align*}
\iint_{\frac{x^2}{a^2}+\frac{y^2}{b^2}+\frac{z^2}{c^2}=1}\frac{1}{\sqrt{\frac{x^2}{a^4}+\frac{y^2}{b^4}+\frac{z^2}{c^4}}}\mathrm{d}S=4\pi abc.
\end{align*}
\end{proof}

\begin{example}
计算
\begin{align*}
\iiint_{\Omega}\frac{1}{(1+x+y+z)^3}\mathrm{d}x\mathrm{d}y\mathrm{d}z,
\end{align*}
这里
\begin{align*}
\Omega=\{(x,y,z)\in\mathbb{R}^3:x,y,z\geqslant0,x+y+z\leqslant1\}.
\end{align*}
\end{example}
\begin{proof}

利用\hyperref[theorem:余面积公式]{余面积公式}，我们有
\begin{align*}
&\iiint_{\Omega}\frac{1}{(1+x+y+z)^3}\mathrm{d}x\mathrm{d}y\mathrm{d}z
=\iiint_{x+y+z\leqslant1}\frac{\chi_{[0,+\infty)^3}(x,y,z)}{(1+x+y+z)^3}\mathrm{d}x\mathrm{d}y\mathrm{d}z \\
&\xlongequal{\text{\hyperref[theorem:余面积公式]{余面积公式}}}\frac{1}{\sqrt{3}}\int_{-\infty}^{1}\mathrm{d}r\iint_{x+y+z=r}\frac{\chi_{[0,+\infty)^3}(x,y,z)}{(1+x+y+z)^3}\mathrm{d}S
=\frac{1}{\sqrt{3}}\int_{0}^{1}\mathrm{d}r\iint_{x+y+z=r}\frac{\chi_{[0,+\infty)^3}(x,y,z)}{(1+r)^3}\mathrm{d}S \\
&=\frac{1}{\sqrt{3}}\int_{0}^{1}\frac{1}{(1+r)^3}\mathrm{d}r\iint_{\substack{x+y+z=r,\\x,y,z\geqslant0}}1\mathrm{d}S
=\frac12\int_{0}^{1}\frac{r^2}{(1+r)^3}\mathrm{d}r=\frac{\ln2}{2}-\frac{5}{16}.
\end{align*}
\end{proof}

\begin{corollary}[Poisson公式]\label{corollary:多元函数Poisson公式}
\begin{enumerate}[(1)]
\item 设$r>0$，$a,b,c\in\mathbb{R}$，且下述积分有意义，则有
\begin{gather*}
\iint_{x^2+y^2\leqslant r^2} f(ax+by)\,g(x^2+y^2)\mathrm{d}x\mathrm{d}y
=\iint_{x^2+y^2\leqslant r^2} f\bigl(\sqrt{a^2+b^2}\,y\bigr)\,g(x^2+y^2)\mathrm{d}x\mathrm{d}y,\\
\oint_{x^2+y^2=r^2} f(ax+by)\,g(x^2+y^2)\mathrm{d}s
=\oint_{x^2+y^2=r^2} f\bigl(\sqrt{a^2+b^2}\,y\bigr)\,g(x^2+y^2)\mathrm{d}s.
\end{gather*}

\item 设$r>0$，$a,b,c\in\mathbb{R}$，且下述积分有意义，则有
\begin{align*}
&\iiint_{x^2+y^2+z^2\leqslant r^2} f(ax+by+cz)\,g(x^2+y^2+z^2)\mathrm{d}x\mathrm{d}y\mathrm{d}z\\
&=\iiint_{x^2+y^2+z^2\leqslant r^2} f\left(\sqrt{a^2+b^2+c^2}\,z\right)g(x^2+y^2+z^2)\mathrm{d}x\mathrm{d}y\mathrm{d}z,
\end{align*}
\begin{align*}
&\iint_{x^2+y^2+z^2=r^2} f(ax+by+cz)\,g(x^2+y^2+z^2)\mathrm{d}S
\\
&=\iint_{x^2+y^2+z^2=r^2} f\left(\sqrt{a^2+b^2+c^2}\,z\right)g(x^2+y^2+z^2)\mathrm{d}S.
\end{align*}
\end{enumerate}
\end{corollary}
\begin{remark}
Poisson公式本质上依赖于换元法，所以实际运用中要注意变通，并不只适用于上述形式。此外,Poisson公式对$n$重积分和$n$重曲面积分也是正确的。
\end{remark}

\begin{example}
设$a,b,c\in\mathbb{R}\setminus\{0\}$，$r>0$，计算
\begin{align*}
\iiint_{x^2+y^2+z^2\leqslant r^2} \bigl(x^2+y^2+z^2\bigr)^p \bigl|ax+by+cz\bigr|^q \mathrm{d}x\mathrm{d}y\mathrm{d}z.
\end{align*}
\end{example}
\begin{remark}
如何判断广义重积分的收敛性？记住固定的操作：加绝对值之后无脑算，算出来收敛就收敛，算出来发散就发散。
\end{remark}
\begin{proof}

直接计算有
\begin{align*}
&\iiint_{x^2+y^2+z^2\leqslant r^2} \bigl(x^2+y^2+z^2\bigr)^p \bigl|ax+by+cz\bigr|^q \mathrm{d}x\mathrm{d}y\mathrm{d}z
\xlongequal{\text{{\hyperref[theorem:余面积公式]{余面积公式}}}}\int_{0}^{r} t^{2p}\mathrm{d}t \iint_{x^2+y^2+z^2=t^2} |ax+by+cz|^q \mathrm{d}S \\
&\xlongequal{\text{\hyperref[corollary:多元函数Poisson公式]{Poisson公式}}}\bigl(a^2+b^2+c^2\bigr)^{\frac q2}\int_{0}^{r} t^{2p}\mathrm{d}t \iint_{x^2+y^2+z^2=t^2} |z|^q \mathrm{d}S 
=2\bigl(a^2+b^2+c^2\bigr)^{\frac q2}\int_{0}^{r} t^{2p}\mathrm{d}t \iint_{\substack{x^2+y^2+z^2=t^2,
z\geqslant0}} z^q \mathrm{d}S \\
&\xlongequal[x=\frac{x^{\prime}}{t},y=\frac{y^{\prime}}{t},z=\frac{z^{\prime}}{t}]{\text{\hyperref[theorem:第一类曲线曲面积分换元法]{第一类曲面积分换元法}}}2\bigl( a^2+b^2+c^2 \bigr) ^{\frac{q}{2}}\int_0^r{t^{2p}\mathrm{d}t\iint_{{x^{\prime}}^2+{y^{\prime}}^2+{z^{\prime}}^2=1,z^{\prime}\geqslant 0}{\frac{{z^{\prime}}^q}{t^q}\cdot \frac{1}{t^2}\mathrm{d}S}}\\
&=2\bigl( a^2+b^2+c^2 \bigr) ^{\frac{q}{2}}\int_0^r{t^{2p+q+2}\mathrm{d}t\iint_{x^2+y^2+z^2=1,z\geqslant 0}{z^q\mathrm{d}S}}\\
&=2\bigl( a^2+b^2+c^2 \bigr) ^{\frac{q}{2}}\int_0^r{t^{2p+q+2}\mathrm{d}t\iint_{x^2+y^2\leqslant 1}{\left( 1-x^2-y^2 \right) ^q\mathrm{d}x\mathrm{d}y}}\\
&\xlongequal{\text{\hyperref[theorem:余面积公式]{余面积公式}}}2\bigl(a^2+b^2+c^2\bigr)^{\frac q2}\int_{0}^{r} t^{2p+q+2}\mathrm{d}t \int_{0}^{1} \mathrm{d}v \int_{x^2+y^2=v^2} \bigl(1-x^2-y^2\bigr)^{\frac{q-1}{2}} \mathrm{d}s \\
&\xlongequal{\text{\hyperref[theorem:第一类曲线积分计算公式]{第一类曲线积分计算公式}}}4\pi\bigl(a^2+b^2+c^2\bigr)^{\frac q2}\int_{0}^{r} t^{2p+q+2}\mathrm{d}t \int_{0}^{1} v\bigl(1-v^2\bigr)^{\frac{q-1}{2}} \mathrm{d}v \\
&=2\pi \bigl( a^2+b^2+c^2 \bigr) ^{\frac{q}{2}}\int_0^r{t^{2p+q+2}\mathrm{d}t\int_0^1{\bigl( 1-v \bigr) ^{\frac{q-1}{2}}\mathrm{d}v}}.
\end{align*}
因此积分收敛等价于
\begin{align}
2p+q+2>-1,\quad \frac{q-1}{2}>-1.\label{eq:conv207}
\end{align}
于是当$p,q$满足\eqref{eq:conv207}时，就有
\begin{align*}
\iiint_{x^2+y^2+z^2\leqslant r^2} \bigl(x^2+y^2+z^2\bigr)^p \bigl|ax+by+cz\bigr|^q \mathrm{d}x\mathrm{d}y\mathrm{d}z
=4\pi \bigl(a^2+b^2+c^2\bigr)^{\frac q2}\frac{(q+1)\,r^{2p+q+3}}{2p+q+3}.
\end{align*}
\end{proof}

\begin{example}
计算
\begin{align*}
\int_{0}^{2\pi}\mathrm{d}\psi \int_{0}^{\pi} e^{\sin\theta(\cos\psi-\sin\psi)}\sin\theta\,\mathrm{d}\theta.
\end{align*}
\end{example}
\begin{proof}

运用\hyperref[theorem:第一类曲面积分的参数方程法]{曲面积分的参数方程}和\hyperref[corollary:多元函数Poisson公式]{Poisson公式}，我们有
\begin{align*}
&\int_0^{2\pi}{\int_0^{\pi}{e^{\sin \theta (\cos \psi -\sin \psi )}\sin \theta \,\mathrm{d}\theta \mathrm{d}\psi}}\xlongequal{\text{\hyperref[theorem:第一类曲面积分的参数方程法]{曲面积分的参数方程}}}\iint_{x^2+y^2+z^2=1}{e^{x-y}\mathrm{d}S}
\\
&\xlongequal{\text{\hyperref[corollary:多元函数Poisson公式]{Poisson公式}}}\iint_{x^2+y^2+z^2=1}{e^{\sqrt{2}\,z}\mathrm{d}S=\int_0^{2\pi}{\int_0^{\pi}{e^{\sqrt{2}\cos \psi}\sin \psi \,\mathrm{d}\psi \mathrm{d}\theta}}}=2\sqrt{2}\,\pi \sinh \sqrt{2}.
\end{align*}
\end{proof}

\begin{theorem}[旋转体表面积和体积公式]\label{theorem:旋转体表面积和体积公式}
\begin{enumerate}[(1)]
\item 穿过直线$l$的平面曲线$L$绕直线$l$旋转一周形成旋转体表面积为
\begin{align*}
\int_L 2\pi d(x,y)\,\mathrm{d}s,
\end{align*}
其中$d$为曲线上的点到直线的距离.

\item 不穿过直线$l$的平面区域$D$绕直线$l$旋转一周形成旋转体体积为
\begin{align*}
\iint_D 2\pi d(x,y)\,\mathrm{d}x\mathrm{d}y,
\end{align*}
其中$d$为区域上的点到直线的距离.
\end{enumerate}
\end{theorem}

\begin{example}
求$L:y=\frac{1}{3}x^3+2x,\,0\leqslant x\leqslant 1$绕$y=\frac{4}{3}x$旋转一周形成旋转体表面积。
\end{example}
\begin{proof}

直接套用\hyperref[theorem:旋转体表面积和体积公式]{旋转体表面积和体积公式}有
\begin{align*}
S &= 2\pi \int_L d(x,y)\mathrm{d}s = 2\pi \int_L \frac{\left|y-\frac{4}{3}x\right|}{\sqrt{1+\left(\frac{4}{3}\right)^2}}\mathrm{d}s = 2\pi \int_{0}^{1} \frac{\left|y-\frac{4}{3}x\right|}{\sqrt{1+\left(\frac{4}{3}\right)^2}}\sqrt{1+(x^2+2)^2}\mathrm{d}x \\
&= 2\pi \int_{0}^{1} \frac{\left|\frac{1}{3}x^3+2x-\frac{4}{3}x\right|}{\sqrt{1+\left(\frac{4}{3}\right)^2}}\sqrt{1+(x^2+2)^2}\mathrm{d}x = \frac{2\pi\sqrt{10}-\pi\sqrt{5}}{3}.
\end{align*}
\end{proof}

\begin{example}
求$y=x^2,\,y=mx\,(m>0)$在第一象限围成区域$D$绕$y=mx$旋转一周形成旋转体体积。
\end{example}
\begin{proof}

直接套用\hyperref[theorem:旋转体表面积和体积公式]{旋转体表面积和体积公式}有
\begin{align*}
V &= 2\pi \iint_D d(x,y)\mathrm{d}x\mathrm{d}y = 2\pi \iint_D \frac{mx-y}{\sqrt{1+m^2}}\mathrm{d}x\mathrm{d}y = 2\pi \int_{0}^{m}\mathrm{d}x \int_{x^2}^{mx} \frac{mx-y}{\sqrt{1+m^2}}\mathrm{d}y
\\
&= \frac{2\pi}{2\sqrt{m^2+1}} \int_{0}^{m} x^2(x-m)^2 \mathrm{d}x 
= \frac{\pi m^5}{30\sqrt{m^2+1}}.
\end{align*}
\end{proof}

\begin{example}
求$(x-2)^2+y^2 \leqslant 1$绕$y$轴旋转一周所得旋转体表面积和体积。
\end{example}
\begin{proof}

直接套用\hyperref[theorem:旋转体表面积和体积公式]{旋转体表面积和体积公式}有
\begin{align*}
S = 2\pi \int_{(x-2)^2+y^2=1} x\mathrm{d}s = 2\pi \int_{(x-2)^2+y^2=1} (x-2+2)\mathrm{d}s 
=4\pi \int_{(x-2)^2+y^2=1}1\,\mathrm{d}s=8\pi^2;
\end{align*}
以及
\begin{align*}
V = 2\pi \iint_{(x-2)^2+y^2\leqslant1}x\,\mathrm{d}x\mathrm{d}y
=2\pi \iint_{(x-2)^2+y^2\leqslant1}(x-2+2)\,\mathrm{d}x\mathrm{d}y 
=4\pi \iint_{(x-2)^2+y^2\leqslant1}1\,\mathrm{d}x\mathrm{d}y=4\pi^2.
\end{align*}
\end{proof}

\begin{example}
设$a>0$，求$\begin{cases}x=a(t-\sin t)\\y=a(1-\cos t)\end{cases},\,t\in[0,2\pi]$和$x$轴围成区域绕直线$y=2a$旋转一周形成的旋转体体积。
\end{example}
\begin{proof}

直接套用\hyperref[theorem:旋转体表面积和体积公式]{旋转体表面积和体积公式}有
\begin{align*}
V &= 2\pi \iint |y-2a|\mathrm{d}x\mathrm{d}y = 2\pi \int_{0}^{2\pi}\mathrm{d}x \int_{0}^{y(x)} (2a-y)\mathrm{d}y \\
&= 2\pi \int_{0}^{2\pi}\left[2a y(x)-\frac{1}{2}y^2(x)\right]\mathrm{d}x
=2\pi a^3 \int_{0}^{2\pi}\left[2(1-\cos t)-\frac{1}{2}(1-\cos t)^2\right]\mathrm{d}(t-\sin t) \\
&=2\pi a^3 \int_{0}^{2\pi}\frac{3-5\cos t+\cos^2 t+\cos^3 t}{2}\mathrm{d}t
=2\pi a^3\cdot \frac{7}{2}\pi=7\pi^2 a^3.
\end{align*}
\end{proof}








\end{document}